% Generated draft: requires mathematical and source review.
\documentclass[11pt]{article}
\usepackage[margin=25mm]{geometry}
\usepackage{fontspec}
\setmainfont{Latin Modern Roman}
\usepackage{amsmath,amssymb,mathtools,unicode-math}
\setmathfont{Latin Modern Math}
\usepackage{enumitem,needspace,xurl,hyperref}
\setlength{\parindent}{0pt}
\setlength{\parskip}{5pt}
\setlength{\emergencystretch}{3em}
\newcommand{\N}{\mathbb{N}}
\newcommand{\Z}{\mathbb{Z}}
\newcommand{\Q}{\mathbb{Q}}
\newcommand{\R}{\mathbb{R}}
\newcommand{\C}{\mathbb{C}}
\newcommand{\F}{\mathbb{F}}
\newcommand{\A}{\mathbb{A}}
\newcommand{\PP}{\mathbb{P}}
\newcommand{\E}{\mathbb{E}}
\newcommand{\eps}{\varepsilon}
\newcommand{\dd}{\,\mathrm{d}}
\begin{document}
\section*{UnsolvedMath ID: 20002003}
\textbf{Source problem code:} catalogue entry
\subsection{Definitions and mathematical statement}
Lauren Williams's source asks how tropicalization of a totally positive variety relates to its cluster algebra structure. A positive part requires a specified positive atlas or coordinate convention; for a Grassmannian it means that the chosen Pl\"ucker coordinates are positive. Its positive tropicalization is obtained from valuations of points with positive leading coefficients over an appropriate real valued extension, with closure and lineality conventions specified. The question concerns both the tropical set and the polyhedral fan chosen on it, which are different data.

The source cites joint work with David Speyer on \(\operatorname{Gr}(k,n)\), relating the \(k=2\) fans to type \(A\) associahedra and the \((3,6),(3,7)\) fans to types \(D_4,E_6\). It relates these to Joshua Scott's cluster structures of the corresponding types. The aid refines the historical statement: after lineality, \(\operatorname{Gr}(2,n)\) gives the Stanley--Pitman fan, combinatorially type \(A\); the positive Pl\"ucker-coordinate fans for \(\operatorname{Gr}(3,6)\) and \(\operatorname{Gr}(3,7)\) are coarsenings, refined using non-Pl\"ucker cluster variables to the \(D_4,E_6\) cluster fans. Combinatorial type is not automatically equality with a particular embedded g-vector fan. Full references and coordinate identifications are not supplied.

The aid's finite criterion is conditional on a known full fan. In a fixed positive Laurent chart, let
\[
 f_i(x)=\sum_{a\in A_i}c_{i,a}x^a,\quad c_{i,a}>0,\qquad
 \tau_i(w)=\min_{a\in A_i}\langle a,w\rangle,
 \quad1\leq i\leq m.
\]
Assume coefficient valuation zero, finite supports, a fixed quotient by common lineality, and that the common fan of linearity \(\mathcal F\) of all \(\tau_i\) is already known to equal the desired cluster g-vector fan in those coordinates. For a retained index subset \(S\subset\{1,\ldots,m\}\), let \(\mathcal F_S\) be the common linearity fan of \(\tau_i\), \(i\in S\); it coarsens \(\mathcal F\). Then
\[
 \mathcal F_S=\mathcal F
\]
if and only if for every pair of adjacent maximal chambers \(C,C'\) of \(\mathcal F\), sharing a facet, at least one retained \(\tau_i\) has different linear gradients on \(C\) and \(C'\). Equivalently, after identifying duplicate exponents and the lineality quotient, some retained coordinate changes its tropical minimizing exponent across every wall.

Failure has a finite local certificate: adjacent chambers on which every retained coordinate has the same linear formula. The result is a test inside a known full positive coordinate presentation; it does not establish the initial cluster-fan identification, produce a cluster structure for an arbitrary positive variety, or imply that Pl\"ucker coordinates alone detect every cluster wall.
\subsection{Short English statement}
The historical Grassmannian examples require distinguishing positive-coordinate fans from their cluster refinements. The aid gives a finite test: a coordinate subfamily preserves a known full cluster fan precisely when it detects every adjacent wall. The initial cluster identification remains an explicit hypothesis.
\subsection{Sources}
\begin{itemize}
\item Local UnsolvedMath registry, permanent ID \(20002003\), source code \texttt{catalogue entry}; canonical collection \texttt{aim-geometry-notes.json}, notes[340].
\item AIM workshop ``Amoebas and tropical geometry,'' Section A.21, ``Positive tropical varieties and cluster algebras,'' contributed by Lauren Williams: \url{https://aimath.org/WWN/amoebas/amoebas.pdf}. The supplied verification reports comparison of the AIM HTML page, PDF and stored record; no external source or PDF was opened in this session.
\item Source attributions: David Speyer and Lauren Williams for positive Grassmannian tropicalizations, Joshua Scott for the stated cluster types. The aid names the Stanley--Pitman fan. No full authoritative titles or bibliography are supplied; original remarks and literature fields are empty, with no proposer-field value or year.
\item Supplied research attempt \(1\), \texttt{partial\_result}, run status \texttt{partially\_solved}, rigor label ``proved,'' novelty confidence ``low.'' Registry classification \texttt{PARTIAL-PROGRESS}, updated 14 August 2026.
\end{itemize}
\end{document}
